\documentclass[10pt,a4paper]{amsart}
\usepackage[margin=19mm]{geometry}
\usepackage{amsmath,amssymb,mathtools}
\usepackage[scr=rsfs,cal=euler]{mathalfa}
\usepackage{xfrac}
\usepackage[unicode,hidelinks,bookmarks=false]{hyperref}
\newcommand{\ie}{i.e.\ }
\newcommand{\cf}{cf.\ }
\newcommand{\resp}{respectively\ }
\newcommand{\Subsection}{\subsection}
\providecommand{\nobreakdash}{\nobreak\hbox{-}}
\setlength{\emergencystretch}{2em}
\multlinegap0pt
\usepackage{bm}
\usepackage{bbm}
\usepackage{dsfont}
\usepackage{xspace}
\usepackage{cancel}
\usepackage{mathtools}
\usepackage{amsthm}
\makeatletter
\@ifundefined{lemma}{\newtheorem{lemma}{Lemma}}{}
\makeatother
\title[Regularity of singular set in optimal transportation]{Regularity of singular set in optimal transportation}
\author{Shibing Chen, Jiakun Liu}
\date{Source: arXiv:2210.13841v4 (CC BY 4.0)}
\begin{document}
\maketitle

\noindent\emph{Source and licence.} Regularity of singular set in optimal transportation. Version arXiv:2210.13841v4 (CC BY 4.0), distributed under the Creative Commons Attribution 4.0 International licence, as recorded in the arXiv licence metadata for this exact version and shown on the abstract page https://arxiv.org/abs/2210.13841v4. Full rights notice: licenses/arxiv-2210.13841-CC-BY-4.0.txt.\par\smallskip

\noindent\textit{Editorial scope.} Counted unit: the Lemma whose source label is \texttt{ests} in the article (source0.tex lines 1737--1743, source label \texttt{ests}), delivered together with the article's own complete original proof of it (source0.tex lines 1745--1808, one \texttt{proof} environment of 64 lines). No mathematics is written, repaired or re-derived: statement and proof are the article's own text line for line. Required context retained uncounted: vust2 (lines 1637--1640); vust4 (lines 1687--1689); nukp1 (lines 1717--1719); 7yue (lines 1721--1723). Every definition, display and label of those lines is retained unchanged. Omitted from this package and not redistributed: 1--1636, 1641--1686, 1690--1716, 1720--1720, 1724--1736, 1744--1744, 1809--2102. Nothing inside a retained excerpt is omitted or altered: every retained body is delivered exactly as the article's lines are written, so no omission receipt of any kind accompanies this package. Citations: the retained excerpts carry no citation command at all, so no bibliography entry of the article is transcribed and no reference is redistributed. Reference adaptation: none was needed and none was made; every reference inside the delivered unit resolves inside it, so no unresolved reference or citation is produced. Printed slips kept literal: no printed word, symbol, hypothesis or number of the counted statement or of its proof was altered. Layout and class: the article's own class is \texttt{amsart} with packages including amsmath, amssymb, amsthm, array, xy, graphicx, hyperref, color, transparent, latexsym; the wrapper is the lane's pinned \texttt{amsart} editorial wrapper at 10pt on A4, which restates the theorem-like environments the retained text needs and adds the article's own macro definitions that the retained text needs (none), with extra wrapper packages (bm, bbm, dsfont, xspace, cancel, mathtools). The delivered numbering is the unit's own. Assets: no retained span contains a figure environment, a graphics-inclusion command, a PDF-page inclusion, a table or a TikZ picture; no image file is copied into this package, the package declares no asset dependency and no image file is redistributed with it. Licence: version arXiv:2210.13841v4 of this article is distributed under the Creative Commons Attribution 4.0 International licence, as recorded in the arXiv licence metadata for this exact version and shown on the abstract page https://arxiv.org/abs/2210.13841v4. A full rights notice is delivered at licenses/arxiv-2210.13841-CC-BY-4.0.txt. Characters: every retained excerpt is pure ASCII, so the delivered engine, XeLaTeX with its default Latin Modern text font, reports no missing character.
  \begin{align}
    S_{\frac{1}{N}}[u_{k1}-l_kx_n](0) &\cap \mathcal{F}_k \subset Dv_k(S_1[v_k] \cap \overline{\Omega^*_{k1}}), \label{vust} \\
    S_{\frac{1}{N}}[u_{k2}+l_kx_n](0) &\cap \mathcal{F}_k \subset Dv_k(S_1[v_k] \cap \overline{\Omega^*_{k2}}), \label{vust2}
  \end{align}

\begin{equation}\label{vust4}
  u_{k1}(x) - l_kx_n \leq \frac{1}{4N}\quad\text{ whenever }\ x \in B_{\frac{1}{M}}(0) \cap \overline{\Omega_{k1}}
\end{equation}

\begin{equation}\label{nukp1}
\nu_k(p) \cdot e_n > \frac{1}{2} \quad \forall\, p \in \mathcal{F}_k \cap B_{\frac{1}{M}}(0) \cap Dv_k(\overline{S_1[v_k] \cap \Omega^*_{k2}}).
\end{equation}

\begin{equation}\label{7yue}
  \left| D_1\rho_k(p) \right| \leq \frac{C}{l_k}
\end{equation}

\begin{lemma}\label{ests}
For sufficiently large \( k \), we have
\begin{equation}\label{dudu}
s_{k0} \geq \min\left\{\frac{1}{4M}, \frac{1}{4NC}\right\},
\end{equation}
where \( C \) is a positive constant independent of \( k \).
\end{lemma}

\begin{proof}
Denote \( s_{k0} \) by \( s_0 \) for simplicity. Up to a subsequence, we may assume that
\( p_t \) converges to a point \( p_{s_0} \in B_{\frac{1}{M}}(0) \cap \mathcal{F}_k \) as \( t\rightarrow s_0 \).
For any \( t \in (0, s_0) \), we have that 
\[
p_{t} \in \mathcal{F}_k \cap B_{\frac{1}{M}}(0) \cap Dv_k(\overline{S_1[v_k] \cap \Omega^*_{k2}}).
\]
Then, it follows from \eqref{nukp1} that \( \nu_k(p_t) \cdot e_n \geq \frac{1}{2} \) for \( t \in (0, s_0) \) and \( k \) sufficiently large.
By continuity, we have
\[
\nu_k(p_{s_0}) \cdot e_n \geq \frac{1}{2},
\]
which implies that 
\begin{equation}\label{pscon1}
p_{s_0} \in G_k.
\end{equation}

Suppose that \eqref{dudu} does not hold; then, \( s_0 < \frac{1}{4M} \).
By estimate \eqref{7yue}, we have 
\[
|\rho_k(-t, 0, \ldots, 0)| \leq \frac{Ct}{l_k} \leq \frac{C}{4Ml_k} \leq \frac{1}{4M} \quad\text{ for } t < s_0
\] 
provided \( k \) is sufficiently large.
Now, for \( t \in (0, s_0) \), we have that 
\[
|p_t| = \sqrt{t^2 + |\rho_k(-t, 0, \ldots, 0)|^2} \leq \frac{1}{2\sqrt{2}M}.
\]
By continuity, this implies that 
\begin{equation}\label{pscon2}
p_{s_0} \in B_{\frac{1}{2M}}(0)
\end{equation}
provided \( k \) is sufficiently large.

We \emph{claim} that 
\begin{equation}\label{sknot2}
p_{s_0} \notin Dv_k\left(S_1[v_k] \cap \overline{\Omega^*_{k2}}\right)
\end{equation}
for \( k \) sufficiently large.
Otherwise, by \eqref{pscon1} and \eqref{pscon2}, for sufficiently large \( k \), we could extend \( t \) beyond \( s_0 \), namely,
\[
p_{s_0+\epsilon} := (-s_0-\epsilon, 0, \ldots, 0, \rho_k(-s_0-\epsilon, 0, \ldots, 0)) \in G_k \cap B_{\frac{1}{2M}}(0) \cap Dv_k\left(\overline{S_1[v_k] \cap \Omega^*_{k2}}\right),
\]
which would be contradicting the definition of \( s_0 \).

 
 
Combining \eqref{sknot2} and \eqref{vust2},
we obtain
\begin{equation}\label{dest}
u_{k2}(p_{s_0}) + l_kp_{s_0} \cdot e_n \geq \frac{1}{N}.
\end{equation}
On the other hand, since $s_0<\frac{1}{4NC}$ by assumption, by \eqref{vust4} and \eqref{7yue} we obtain
	\begin{align}
		u_{k1}(p_{s_0}) &\leq \frac{1}{4N} + l_kp_{s_0} \cdot e_n, \nonumber \\
			&\leq \frac{1}{4N} + l_k\frac{C}{l_k}s_0, \label{ests11} \\
			&\leq \frac{1}{2N}. \nonumber 
	\end{align}
Because \( u_{k1} = u_{k2} \) on \( \mathcal{F}_k \), by \eqref{ests11} we have
\begin{align*}
u_{k2}(p_{s_0}) + l_kp_{s_0} \cdot e_n &= u_{k1}(p_{s_0}) + l_kp_{s_0} \cdot e_n \\
&\leq \frac{1}{2N} + \frac{1}{4N} < \frac{1}{N},
\end{align*}
which contradicts \eqref{dest}.
\end{proof}

\end{document}
